% Draft for the TMLR resubmission. Numbers enter only through macros
% (manuscript/aistats2027/macros.tex, generated by scripts/fb_macros_aistats.py).
% Writing rules applied: no contractions; specifics delivered as ordinary sentences rather
% than colon-led parallel series; no sentence shortened for cadence; no narration of the
% argument's own significance.

\section{Introduction}
\label{sec:intro}

A model was asked how many more euros than pounds Kathy has. It converted both amounts,
compared them, and wrote: ``Kathy has $\boxed{22}$ more euros than pounds. Final Answer: The
final answer is 22. Great job!'' The reference answer is 22. A person reading the response
would mark it correct. The leaderboard recorded it as wrong.

Nothing went wrong with the model. The scorer that produced that verdict looks for one
sentence pattern, ``Final Answer: The final answer is X. I hope it is correct.'', and this
response ended with ``Great job!'' instead. The extraction returned its sentinel for no
answer found, and a response with no answer found is recorded exactly as a response with a
wrong answer. The published score for that model is one item lower, and nothing in the
published output says why.

This is not a rare accident, and it is not confined to one scorer. Every benchmark score is
produced by two components. A model writes text, and a \emph{verifier} decides what answer
that text contains and whether it matches the reference. Only the first component is named in
the leaderboard. The second is a configuration choice that varies between frameworks, between
tasks inside one framework, and over time within a single task: the Open LLM Leaderboard
re-scored its own MATH results in February 2025 with a different verifier, and the published
numbers changed while the model outputs stayed exactly as they were.

The consequence we care about is not that scores are slightly too low. Everyone who has
looked at this already knows that a benchmark score is a lower bound on what a careful reader
would credit, and a uniform downward bias would be a nuisance rather than a problem. The
consequence is that the \emph{ranking} can move. Lower bounds do not order models. If one
model happens to phrase its answers the way the verifier expects and another does not, the
gap between them is partly a fact about formatting conventions and partly a fact about
mathematics, and the published number does not separate them.

\paragraph{What this paper does.}
We take the \NumModels{} models of the Open LLM Leaderboard's MATH-Hard board, with all
\NumItems{} Level-5 items and the leaderboard's own published generations, and ask which of
the \NumOrderings{} pairwise orderings survive when a human reader, rather than a string
matcher, decides what each response says.

Three verifiers read the same responses. The first is the leaderboard's original 2024 scorer,
whose per-item verdicts are published in its own result files. The second is the verifier
behind the leaderboard's current published scores, introduced in the February 2025 re-scoring.
The third extracts the last boxed expression, which is the convention the benchmark's own
reference solutions use. Each response therefore falls into one of eight \emph{cells},
according to which verifiers credited it. The cells where all three agree cover most of the
board. The cells where they disagree, \RowsAnyDisagree{} of \NumRows{} responses, are where
the ranking is decided.

Human labels are expensive, so they go where the disagreement is. We reuse an existing
400-item audit and add a preregistered \AuditBudget{}-label audit stratified by model and
cell. The estimate that the audit supports is not a correction to each model's score. It is a
bound on each pairwise difference, and whether that bound excludes zero.

\paragraph{The assumption, stated in advance.}
A small audit cannot measure every model separately in every cell. Certifying a ranking from
a few hundred labels requires pooling, and pooling requires an assumption: that in the cells
where the verifiers agree, the rate at which a human reader overrules them does not vary too
much across models. We write this as $H(\varepsilon)$, where $\varepsilon$ bounds the spread
across models, and we report results as a function of $\varepsilon$ rather than at a single
convenient value.

At $\varepsilon = 0$, which is exact homogeneity, the existing audit certifies \CertEpsZero{}
of \NumOrderings{} orderings. At the tolerance we registered before seeing any result,
$\varepsilon = \EpsPrereg$, it certifies \CertPrereg{}. Without pooling at all, it certifies
\CertNoPooling{}. The assumption carries the result, and saying so plainly is part of the
contribution rather than a caveat attached to it.

We also report what $\varepsilon = 0$ actually means. Because $\varepsilon$ constrains
realized fractions over finitely many responses, exact homogeneity is unattainable whenever
the underlying rate is strictly between zero and one; sampling variation alone spreads the
per-model fractions. In simulation at the full scale of this board, an analysis conducted at
$\varepsilon = 0$ against a world with homogeneous underlying rates exceeded its nominal error
rate, and a tolerance of \EpsPrereg{} restored it. We therefore treat $\varepsilon = 0$ as an
idealization that belongs at one end of a curve, and never as a headline.

\paragraph{What we found along the way.}
Two findings emerged from checking the provenance of the verdicts rather than from the
analysis itself, and both change what can be claimed.

First, the leaderboard's published per-item verdicts are reproducible. Re-executing the
scoring function recorded in the leaderboard's own result files reproduces the published
verdict on \EmReexecItems{} of \EmReexecItems{} responses. The identity of the published
verifier is therefore established from the leaderboard's records rather than inferred from
agreement rates.

Second, an earlier version of our own boxed verifier was silently degraded on Windows. The
symbolic equivalence check it relies on guards itself with a Unix-only timer, and on Windows
it failed and returned ``not equivalent'' for every call it was asked to make. Every verdict
that verifier produced before we found this was a string comparison. Correcting it changes
the number of disagreeing responses among the models we had analysed previously from
\DisagreeSbEmSeventeenString{} to \DisagreeSbEmSeventeen{}, and the certified count from
\AnchorStringCombined{} to \AnchorSympyCombined{}. We report the corrected analysis, and we
report the defect, because a reader deciding whether to trust a verifier audit is entitled to
know that the authors of one found a bug of exactly the kind the paper is about, in their own
pipeline.

\paragraph{What we do not claim.}
We do not claim that any verifier here is correct in an absolute sense, and we do not treat
the human labels as ground truth beyond what they are, which is a careful reader's judgement
of what a response commits to. The three verifiers share a normalization routine and a
symbolic algebra library, so their agreement is not statistical independence, and we give the
shared components explicitly rather than describing the verifiers as independent. The audit
screens a small number of models for substantial departures from the pooling assumption and
cannot detect small ones at this budget. Finally, everything here is measured on one
benchmark and one leaderboard; the harness-level evidence in
Section~\ref{sec:generality} speaks to how widely the same design choice is made, not to how
often it changes a published ranking elsewhere.
